% ---------------------------------------------------------------------------
%  QML Lab research project 13: Learning from Classical Shadows
%  Build:  latexmk -pdf main.tex   (needs qmlproject.sty, harvardqml.bst, references.bib)
% ---------------------------------------------------------------------------
\documentclass[11pt,a4paper]{article}
\usepackage{qmlproject}
\qmlheader{Project 13: Learning from Classical Shadows}

\hypersetup{pdftitle={QML Lab project 13: Learning from Classical Shadows}}

\begin{document}

\begin{center}
  {\small\textcolor{black!55}{QML Lab \textbullet\ research project 13}}\\[0.5em]
  {\LARGE\bfseries Learning from Classical Shadows}\\[0.6em]
  {\large Predicting many properties of a quantum state from few randomised measurements.}\\[0.6em]
  {\small Living document \textbullet\ last updated \today}
\end{center}

\begin{projectcard}
\small
\begin{tabularx}{\linewidth}{@{}>{\bfseries}l X >{\bfseries}l X@{}}
Quadrant  & \QC\ (quantum data or quantum algorithm, classical learner) & Track & Quantum data \\
Difficulty & Hard & Tools & PennyLane, NumPy, scikit-learn \\
Builds on & Modules 2 and 8 of the QML Lab & Workspace & \path{qml_projects/classical-shadows/} \\
\end{tabularx}
\end{projectcard}

\paragraph{Summary.}
Full quantum state tomography needs a number of measurements that grows exponentially with the number
of qubits. Classical shadows replace it with randomised single-qubit measurements and simple classical
post-processing, and predict many properties at once with far fewer measurements. This project
measures how estimation error scales, compares shadows with tomography, and uses shadow-based features
to train a classical phase classifier, linking it to the QCNN of Project 12.

% ===========================================================================
\section{Background and key concepts}
% ===========================================================================
\subsection{Randomised measurements}
Randomised measurement protocols apply random unitaries before measuring and learn from the
statistics of the outcomes \citep{elben2023}. In the classical-shadow protocol of
\citet{huang2020shadows}, each round applies a random single-qubit Clifford $U_j$ to every qubit (a random
choice of the $X$, $Y$ or $Z$ basis), measures, and stores the snapshot
\begin{equation}
  \hat\rho=\bigotimes_{j=1}^{n}\Big(3\,U_j^\dagger\ket{\hat b_j}\!\bra{\hat b_j}U_j-I\Big),
\end{equation}
which is an unbiased estimator of $\rho$: $\E[\hat\rho]=\rho$.

\subsection{Sample complexity}
Expectation values $o_i=\Tr(O_i\rho)$ are estimated by median-of-means over snapshots. To predict $M$
observables to accuracy $\epsilon$ with high probability, it suffices to take
\begin{equation}
  N_{\text{snap}}=O\!\Big(\frac{\log M}{\epsilon^2}\max_i\lVert O_i\rVert_{\text{shadow}}^2\Big),
  \qquad \lVert O\rVert_{\text{shadow}}^2\le 3^k\lVert O\rVert_\infty^2
\end{equation}
for $k$-local observables with Pauli measurements. Local properties are cheap; global ones are not.
Non-linear quantities such as the purity $\Tr\rho_A^2$ and the second R\'enyi entropy
$S_2(\rho_A)=-\log\Tr\rho_A^2$ follow from pairs of independent snapshots.

\subsection{Machine learning on shadows}
\citet{huang2022science377} proved that classical machine learning trained on classical-shadow data
can efficiently predict ground-state properties and classify phases of matter. Neural-network
tomography \citep{torlai2018} is an alternative that learns a model of the state from measurements.
Both are \QC: the data comes from quantum states, the learning is classical.

% ===========================================================================
\section{Research questions}
% ===========================================================================
The project is open-ended: these questions are the starting point, not the finish line. Add new
ones to the README and the journal as they come up.

\begin{rqlist}
  \item How does the estimation error of local observables and of subsystem entropy scale with the number of snapshots?
  \item Which observables are cheap to estimate and which need prohibitively many snapshots?
  \item Can a classical model trained on shadow features classify Ising phases as well as the QCNN of Project 12, for the same number of state copies?
\end{rqlist}

% ===========================================================================
\section{Method}
% ===========================================================================
\subsection{Data}
GHZ states, random product states and Ising ground states (shared with Project 12) on 4 to 10 qubits;
$10^2$ to $10^4$ snapshots per state, generated with PennyLane's classical-shadow measurement.

\subsection{Models}
Shadow estimators for $\langle X_i\rangle$, $\langle Z_iZ_{i+1}\rangle$, the fidelity with a target state
and $S_2$ of half the chain; full tomography by linear inversion for $n\le4$; a logistic regression or
small MLP on shadow-estimated correlations for phase classification.

\subsection{Baselines}
Exact values from the state vector; full tomography with the same number of measurements; direct
measurement in a fixed basis for each observable.

\subsection{Metrics}
\begin{itemize}
  \item Absolute estimation error versus number of snapshots (log--log).
  \item Snapshots needed to reach $\epsilon=0.05$.
  \item Phase classification accuracy versus number of state copies.
\end{itemize}

% ===========================================================================
\section{Planned experiments}
% ===========================================================================
Each experiment becomes one \texttt{[[experiments]]} entry in the workspace, so its runs, parameters
and metrics are stored and compared in the portal.

\begin{table}[h]
\centering\small
\begin{tabularx}{\linewidth}{@{}l>{\raggedright}X>{\raggedright}X>{\raggedright\arraybackslash}X@{}}
\toprule
ID & Experiment & What varies & Main metrics \\
\midrule
  E1 & Local observables & $N_{\text{snap}}=10^2..10^4$ & error scaling \\
  E2 & Entanglement & $S_2$ of half chain & error vs $N_{\text{snap}}$ \\
  E3 & Tomography comparison & $n\le4$, equal budget & error \\
  E4 & Phase classifier & shadow features of Ising states & accuracy vs copies \\
  E5 & Readout noise & bit-flip probability & bias \\
\bottomrule
\end{tabularx}
\caption{Experiment plan. Results go into the journal and the results log (\cref{sec:log}).}
\end{table}

% ===========================================================================
\section{Running it in the QML Lab}
% ===========================================================================
The workspace lives in \path{qml_projects/classical-shadows/}. Scripts import three helpers from
\path{qml_lab/plots.py}:
\begin{itemize}
  \item \texttt{params()} reads the parameters from the portal's run form (defaults when run by hand);
  \item \texttt{save()} stores figures with the run;
  \item \texttt{log\_metrics()} records numbers that appear in the run table.
\end{itemize}

\begin{lstlisting}
qml_projects/classical-shadows/
  project.toml          metadata + experiments
  README.md             goal, open questions, ideas
  experiments/          runnable scripts
  docs/main.tex         this document
\end{lstlisting}

First experiment: This experiment is \emph{planned}: write the script, add the block below to \texttt{project.toml}, then run \texttt{python manage.py sync\_projects}.

\begin{lstlisting}[language={}]
[[experiments]]
slug = "shadow-scaling"
title = "Error scaling of classical-shadow estimates"
script = "experiments/e01_shadow_scaling.py"
figure = "error_scaling.png"

[experiments.params]
n_qubits = 6
n_snapshots = 1000
state = "ising"
field = 1.0
seed = 0
\end{lstlisting}

\begin{lstlisting}[language=bash]
# from the portal folder: pick up changes to project.toml
python manage.py sync_projects
cd qml_projects/classical-shadows
# run by hand with the default parameters
PYTHONPATH=../../qml_lab python experiments/e01_shadow_scaling.py
\end{lstlisting}

% ===========================================================================
\section{Quality bar}
% ===========================================================================
\begin{itemize}
  \item A classical baseline with a similar number of trainable parameters, tuned with the same care \citep{bowles2024}.
  \item At least five random seeds per configuration, reported as mean $\pm$ standard deviation.
  \item How results change with the number of qubits.
  \item At least one run under a realistic noise model.
  \item Fit and report the error-scaling exponent; it should be close to $-1/2$.
\end{itemize}

% ===========================================================================
\section{Open directions}
% ===========================================================================
\begin{itemize}
  \item Derandomised shadows for a fixed set of observables.
  \item Shadows from real hardware, including readout mitigation.
  \item Shared data set with Project 12 for a direct \QC\ versus \QQ\ comparison.
\end{itemize}

% ===========================================================================
\section{Risks and limitations}
% ===========================================================================
\begin{itemize}
  \item Global observables need exponentially many snapshots; choose observables with this in mind.
  \item Readout errors bias estimates unless mitigated.
  \item Median-of-means parameters affect the error; keep them fixed across comparisons.
\end{itemize}

% ===========================================================================
\section{Results log}
\label{sec:log}
% ===========================================================================
Summarise findings here as they become stable; day-to-day notes belong in the portal journal.

\begin{table}[h]
\centering\small
\begin{tabularx}{\linewidth}{@{}l l >{\raggedright}X >{\raggedright\arraybackslash}X@{}}
\toprule
Date & Run & Finding & Next step \\
\midrule
 & & & \\
 & & & \\
 & & & \\
\bottomrule
\end{tabularx}
\end{table}

\bibliographystyle{harvardqml}
\bibliography{references}

\end{document}
